\chapter{Implementación del firmware de bajo nivel en Arduino Nano}
\label{ch:setup_arduino}

Este capítulo describe el firmware \texttt{ROSArduinoBridge} implementado en el \gls{Arduino} Nano para el control de bajo nivel del robot. El firmware cumple cuatro funciones principales: lectura de encoders, control PID de velocidad por rueda, comunicación serial con \gls{ROS2} mediante \texttt{ros2\_control} y auto-detención de seguridad ante pérdida de comandos.

%------------------------------------------------------------------
\section{Arquitectura modular del firmware}

El firmware se organiza en módulos con responsabilidades separadas: \texttt{ROSArduinoBridge.ino} coordina la comunicación serial y el ciclo principal; \texttt{commands.h} define los comandos; \texttt{encoder\_driver} realiza la lectura de encoders; \texttt{motor\_driver} acciona el L298N; y \texttt{diff\_controller.h} implementa el controlador PID. Esta organización facilita la trazabilidad entre el modelo cinemático, el firmware y la interfaz con \texttt{ros2\_control}.

\utemFigurePropia{2.Texto/Imag/diagrama_arquitectura_firmware.png}{0.85\textwidth}
{Arquitectura modular del firmware \texttt{ROSArduinoBridge}.}
{fig:arquitectura_firmware}

%------------------------------------------------------------------
\section{Constantes físicas y conversión de unidades}

El firmware utiliza las mismas constantes físicas definidas para el modelo cinemático y el gemelo digital: \texttt{WHEEL\_DIAMETER} = 0.067~m y \texttt{ENC\_TICKS\_PER\_REV} = 1980. A partir de ellas, se calcula el factor de conversión entre velocidad lineal en m/s y velocidad de rueda en \textit{ticks/frame}:

\begin{align}
\texttt{MPS\_TO\_TICKS\_FACTOR} &=
\frac{\texttt{ENC\_TICKS\_PER\_REV}}
{\pi \cdot \texttt{WHEEL\_DIAMETER} \cdot \texttt{PID\_RATE}}
\notag \\
&=
\frac{1980}{\pi \times 0.067 \times 30}
\approx 313.56
\frac{\text{ticks/frame}}{\text{m/s}}
\label{eq:mps_to_ticks}
\end{align}

El lazo de control opera a 30~Hz, por lo que el período de muestreo es aproximadamente $T_s=1/30$~s.

%------------------------------------------------------------------
\section{Comunicación serial y lectura de encoders}
\label{sec:serial_firmware}

La comunicación entre la Raspberry Pi y el \gls{Arduino} se realiza mediante un protocolo serial textual a 57600~baud. Durante la operación normal, los comandos principales son: lectura de encoders, envío de referencias de velocidad, carga de ganancias PID y aplicación de PWM directo para pruebas en lazo abierto. El protocolo completo se documenta en el \ref{anx:firmware}.

\subsection{Obtención experimental de cuentas por vuelta}
\label{sec:cpr_experimental}
\label{sec:cpr_promedios}

Para obtener una odometría consistente, se midió experimentalmente el número de cuentas por vuelta completa de rueda. Los promedios obtenidos fueron:

\begin{align}
\bar{N}_R &= \frac{1}{10}\sum_{i=1}^{10} n_{R(i)} = 1976
\label{eq:promedio_cuentas_R} \\
\bar{N}_L &= \frac{1}{10}\sum_{i=1}^{10} n_{L(i)} = 1969
\label{eq:promedio_cuentas_L}
\end{align}

Dado que la diferencia entre ruedas es inferior al 0.4\%, se adopta un valor unificado de \texttt{ENC\_TICKS\_PER\_REV} = 1980 cuentas/vuelta. Las muestras experimentales se incluyen en el \ref{anx:firmware}.

%------------------------------------------------------------------
\section{Control PID en el firmware}
\label{sec:motorcontrol_pid}

El controlador de velocidad se ejecuta de forma independiente para cada rueda. La referencia y la medición se expresan en \textit{ticks/frame}. Para cada rueda $j \in \{R,L\}$, el error se define como:

\begin{equation}
e_j[k] = m_j^{ref}[k] - m_j[k],
\qquad
m_j[k] = \texttt{Encoder}_j[k] - \texttt{PrevEnc}_j[k]
\label{eq:error_pid}
\end{equation}

La ley implementada corresponde a la forma incremental con derivada sobre la medición definida en el Marco Teórico. El firmware incluye anti-windup condicional, zona muerta para reducir ciclos límite por cuantización y saturación de la salida PWM. Las ganancias definitivas son:

\begin{table}[H]
\centering
\caption{Ganancias PID definitivas codificadas en el firmware.}
\label{tab:ganancias_pid}
\small
\begin{tabular}{lcc}
\toprule
\textbf{Parámetro} & \textbf{Rueda izquierda} & \textbf{Rueda derecha} \\
\midrule
$K_p$ & 16.00 & 16.42 \\
$K_d$ & 20.30 & 20.05 \\
$K_i$ & 0.075 & 0.001 \\
$K_o$ & 50 & 50 \\
\bottomrule
\end{tabular}
\end{table}

El divisor $K_o$ normaliza la acción de control, por lo que las ganancias efectivas corresponden a $K_p/K_o$, $K_d/K_o$ y $K_i/K_o$. La convención de orden usada por el firmware es $K_p$, $K_d$, $K_i$, $K_o$, heredada de \texttt{ROSArduinoBridge}.

%------------------------------------------------------------------
\section{Sintonización del controlador PID asistida por PSO}
\label{sec:sintonizacion_pso}

Las ganancias PID se determinaron mediante un procedimiento híbrido. Primero se capturó la respuesta al escalón de cada rueda en lazo abierto; luego se identificó un modelo FODT mediante PSO; posteriormente se optimizaron las ganancias del controlador sobre el modelo identificado; finalmente, las ganancias fueron refinadas experimentalmente en el hardware real.

\subsection{Captura de respuesta en lazo abierto}
\label{sec:captura_lazo_abierto}

La captura se realizó aplicando un escalón de PWM sin controlador PID activo. La velocidad lineal de cada rueda se calculó como:

\begin{equation}
v_j[k] =
\frac{\Delta\texttt{enc}_j[k]}{N_e}
\cdot
\frac{\pi d_w}{T_{\mathrm{cap}}},
\qquad j \in \{R,L\}
\label{eq:vel_lazo_abierto}
\end{equation}

donde $T_{\mathrm{cap}}=0.1$~s, $d_w=0.067$~m y $N_e=1980$ cuentas/vuelta. La variable $v_j[k]$ representa velocidad lineal de rueda, no velocidad lineal del robot.

\subsection{Identificación y optimización}

La identificación FODT se realizó mediante integración discreta del modelo:

\begin{equation}
y[k] =
y[k-1] +
\frac{T_{\mathrm{cap}}}{\tau}
\left(K_m u[k-n_d] - y[k-1]\right)
\label{eq:euler_fodt}
\end{equation}

donde $n_d=\lfloor L_d/T_{\mathrm{cap}}\rfloor$ representa el retardo en muestras. Luego, un segundo PSO buscó las ganancias $(K_p,K_d,K_i)$ que minimizan una función de costo basada en error cuadrático medio y sobreimpulso.

\begin{algorithm}[H]
\caption{Sintonización PID asistida por PSO}
\label{alg:sintonizacion}
\begin{algorithmic}[1]
\Require Datos de lazo abierto $(u_{\mathrm{pwm}}[k],v_j^{med}[k])$
\Ensure Ganancias $(K_p,K_d,K_i)$ por rueda
\State Identificar parámetros FODT $(K_m,\tau,L_d)$ mediante PSO
\State Simular el PID del firmware sobre el modelo FODT
\State Evaluar la función de costo con MSE y sobreimpulso
\State Seleccionar las mejores ganancias del enjambre
\State Refinar experimentalmente las ganancias en hardware real
\State \Return $(K_p,K_d,K_i)$
\end{algorithmic}
\end{algorithm}

\subsection{Refinamiento experimental}
\label{subsec:refinamiento_iterativo}

Las ganancias obtenidas por PSO se usaron como punto de partida. Luego se ajustaron sobre el robot real para compensar efectos no modelados, como fricción estática, zona muerta del L298N y diferencias mecánicas entre ruedas. El resultado final fueron ganancias asimétricas, lo que permite compensar mejor el comportamiento individual de cada motor.

%------------------------------------------------------------------
\section{Bucle principal y auto-detención}
\label{sec:bucle_principal}

El bucle principal atiende el parser serial, ejecuta el PID a 30~Hz y evalúa la auto-detención. Si transcurren más de 2000~ms sin recibir comandos de movimiento, el firmware detiene los motores y desactiva la bandera de movimiento. Esta condición evita que el robot continúe desplazándose ante pérdida de conexión o falla de algún nodo de navegación.

%------------------------------------------------------------------
\section{Integración con ROS 2 mediante \texttt{diffdrive\_arduino}}
\label{sec:integracion_ros2}

La integración con \gls{ROS2} se realiza mediante el plugin \texttt{DiffDriveArduino}, usado por \texttt{ros2\_control}. El plugin traduce velocidades angulares de rueda hacia \textit{ticks/frame}, envía dichas referencias al \gls{Arduino} y recibe los conteos de encoder para actualizar posición, velocidad y odometría.

La consistencia entre plugin y firmware depende de tres parámetros compartidos: \texttt{enc\_counts\_per\_rev} = 1980, \texttt{loop\_rate} = 30~Hz y \texttt{baud\_rate} = 57600. Estos valores permiten que la odometría publicada por \texttt{diff\_drive\_controller} sea coherente con la medición de encoders y con el modelo cinemático.

%------------------------------------------------------------------
\section*{Conclusión del capítulo}

El firmware implementado permite leer encoders, ejecutar control PID por rueda, comunicarse con \gls{ROS2} y detener el robot ante pérdida de comandos. La sintonización híbrida PSO más refinamiento experimental permitió obtener ganancias asimétricas adecuadas para el hardware real. Esta capa de bajo nivel constituye la base para la odometría, el control diferencial, el SLAM y la navegación autónoma.